\documentclass[11pt,a4paper,sans,twocolumn]{moderncv}
\moderncvstyle{classic}
\moderncvcolor{blue}
\usepackage[T1]{fontenc}
\usepackage[polish]{babel}
\usepackage{lmodern}
\usepackage[scale=0.83]{geometry}

\firstname{Sherlock}
\familyname{Holmes}
\title{Analityk śledczy | 2026}
\address{Londyn, świat literacki}{}
\email{s.holmes@example.org}

\begin{document}
\makecvtitle

\section{Profil zawodowy}
\cvitem{}{Analityk obserwacji i wnioskowania. Łączę ślady materialne,
zeznania oraz sekwencję zdarzeń w hipotezy, które można sprawdzić.
Ubiegam się o stanowisko analityka śledczego w fikcyjnym Archiwum
Nierozwiązanych Zagadek.}

\section{Doświadczenie i projekty}
\cventry{Opowiadanie}{Sprawa Ligi Rudowłosych}{Doradca śledczy}{Londyn}{}{
  Zestawiłem harmonogram pozornie niewinnej pracy z otoczeniem miejsca
  zdarzenia. Wskazałem hipotezę dotyczącą celu zatrudnienia i
  zweryfikowałem ją w terenie.}
\cventry{Opowiadanie}{Sprawa błękitnego karbunkułu}{Doradca śledczy}{Londyn}{}{
  Odtworzyłem drogę przedmiotu między kolejnymi osobami. Oddzieliłem
  relacje świadków od obserwowalnych faktów i sprawdziłem zgodność
  chronologii.}

\section{Umiejętności}
\cvitem{Analiza}{formułowanie hipotez, porównywanie śladów, kontrola sprzeczności}
\cvitem{Komunikacja}{zwięzły opis przesłanek i stopnia pewności wniosku}
\cvitem{Organizacja}{prowadzenie notatek z obserwacji i porządkowanie osi czasu}
\cvitem{Organizacja}{prowadzenie notatek z obserwacji i porządkowanie osi czasu}
\section{Przygotowanie do roli}
\cvitem{Nauka}{studium metod obserwacyjnych i dokumentowanie przypadków;
  bez przypisywania postaci niepotwierdzonego dyplomu}

\section{Informacja o charakterze CV}
\cvitem{}{Postać i sprawy pochodzą z opowiadań Arthura Conana Doyle'a.
  Stanowisko, profil pracodawcy, adres e-mail i język kompetencji są
  dopowiedzeniem na potrzeby ćwiczenia.}

\end{document}
